% Generated draft: requires mathematical and source review.
\documentclass[11pt]{article}
\usepackage[margin=25mm]{geometry}
\usepackage{fontspec}
\setmainfont{Latin Modern Roman}
\usepackage{amsmath,amssymb,mathtools,unicode-math}
\setmathfont{Latin Modern Math}
\usepackage{enumitem,needspace,xurl,hyperref}
\setlength{\parindent}{0pt}
\setlength{\parskip}{5pt}
\setlength{\emergencystretch}{3em}
\newcommand{\N}{\mathbb{N}}
\newcommand{\Z}{\mathbb{Z}}
\newcommand{\Q}{\mathbb{Q}}
\newcommand{\R}{\mathbb{R}}
\newcommand{\C}{\mathbb{C}}
\newcommand{\F}{\mathbb{F}}
\newcommand{\A}{\mathbb{A}}
\newcommand{\PP}{\mathbb{P}}
\newcommand{\E}{\mathbb{E}}
\newcommand{\eps}{\varepsilon}
\newcommand{\dd}{\,\mathrm{d}}
\begin{document}
\section*{UnsolvedMath ID: 174}
\textbf{Source problem code:} catalogue entry
\subsection{Definitions and mathematical statement}
For an integer $N\geq1$, write $[N]=\{1,\ldots,N\}$. Let $p_1,\ldots,p_{1000}$ be distinct primes with
\[2\leq p_1<p_2<\cdots<p_{1000}<N^{9/10}.\]
If $R_i\subseteq\Z/p_i\Z$ is the set of residue classes removed at the $i$th prime, the surviving set is
\[S=\{n\in[N]: n\bmod p_i\notin R_i\text{ for every }1\leq i\leq1000\}.\]
The source asks whether removing ``half the residue classes'' for each $p_i$ forces
\[|S|\leq\frac{N}{10}.\]
This is not yet a uniquely defined assertion: for odd $p_i$ the number $p_i/2$ is not an integer, so the sizes $|R_i|$ require a rounding or other convention. The intended quantifier over the choices of $R_i$ is also omitted. The notation above treats removed classes as distinct members of a subset; it does not choose their sizes or resolve the quantifier. The source also does not explicitly state whether its conclusion is for every admissible $N$ or only for sufficiently large $N$. No one of these variants is silently asserted here.

The supplied literature triage, checked 2026-08-17, identifies the local code with Green's Problem 44 and says that the same parity and quantifier defects occur in the maintained source. It reports Erd\H{o}s's observation that the intended assertion is affirmative by the large sieve when all 1000 primes are below $N^{1/2}$, and reports no further progress for the $N^{9/10}$ version. These statements concern the intended, convention-dependent variant. The registry is labelled open and the triage partially solved / PARTIAL-PROGRESS; the review assigns formulation-dependent uncertainty to the literal wording.
\subsection{Short English statement}
After sieving $\{1,\ldots,N\}$ by removing about half the residue classes for each of 1000 primes below $N^{9/10}$, must at most one tenth of the integers remain? The rounding convention and quantifiers over the removed classes must first be specified.
\subsection{Sources}
\begin{itemize}
\item Local UnsolvedMath registry, permanent ID 174, source code catalogue entry, ``Sieving by Many Small Primes,'' in Ben Green's 100 Open Problems; supplied statement and background.
\item B. Green, \emph{100 Open Problems}, Problem 44, version retrieved 2026-08-17 according to the supplied maintained-tracker reference: \url{https://people.maths.ox.ac.uk/greenbj/papers/open-problems.pdf}. The record says this source uses the same terse wording and records the large-sieve observation for primes below $N^{1/2}$.
\item Supplied literature triage checked 2026-08-17, PARTIAL-PROGRESS; it flags parity and quantifier defects rather than repairing them, and cautions that dated triage does not exclude later or unindexed results.
\end{itemize}
\end{document}
